\maketitle

\begin{abstract}
Given a measure-preserving system $(X, \mathcal{B}, \mu, T)$ and a bounded
observation function $h \in L^\infty(X, \mu)$, the time series
$h(x), h(Tx), h(T^2 x), \ldots$ is all an observer sees.  The question is
when this is enough: when can the observer recover the full state space $X$
from the time-delayed record?

The answer is a three-way equivalence.  State-space recovery —
$\mathcal{O}_h := \sigma(\{h \circ T^n : n \in \mathbb{Z}\}) = \mathcal{B}$
modulo $\mu$ — holds if and only if the algebra generated by the delay orbit
$\{h \circ T^n\}$ is dense in $L^2(X, \mu)$, if and only if the delay map
$\Phi_h : x \mapsto (h(T^n x))_{n \in \mathbb{Z}}$ is a measure-theoretic
embedding $X \hookrightarrow \mathbb{R}^\mathbb{Z}$.

The equivalence of the first two conditions is the \emph{density bridge}: for
any subalgebra $\mathcal{A} \subseteq L^\infty$ containing the constant $1$,
$\overline{\mathcal{A}}^{L^2} = L^2(X, \sigma(\mathcal{A}), \mu)$.  The proof
is an instance of the functional monotone class theorem.

The cyclic vector condition $\overline{\mathrm{span}}\{U_T^n h\} = L^2$ implies
reconstruction but is strictly stronger: it is a spectral condition on the
Koopman operator $U_T$, while reconstruction is a condition on the Boolean
algebra of observable events.  They coincide precisely when $U_T$ has simple
$L^2$-spectrum.

When reconstruction holds, the compact Stone space of the observable algebra —
introduced in~\cite{mahon_paper1} as a technical device for deriving
$\sigma$-additivity from compactness — is measure-theoretically isomorphic to
$X$ itself.  Compared to the classical Takens embedding theorem~\cite{Takens1981},
the reconstruction theorem requires only measurability and uses all delays
simultaneously; the algebraic density condition replaces the dimension count.
\end{abstract}

\begin{quote}\small
\textbf{2020 Mathematics Subject Classification.}
37A05 (primary); 28A60, 06E15, 37A30, 46E30 (secondary).

\textbf{Keywords.}
reconstruction theorem, delay embedding, observable algebra, density bridge,
measure-theoretic Takens, Stone duality, Koopman operator, measure-preserving system
\end{quote}

% ====================================================================
\section{Introduction}
\label{sec:intro}

An observer of a measure-preserving system $(X, \mathcal{B}, \mu, T)$ has
access to a bounded observation function $h : X \to \mathbb{R}$, and from a
single initial point $x$ sees the time series $h(x), h(Tx), h(T^2x), \ldots$
The prior papers established that coherent structured observations determine a
probability measure and that measure determines the dynamics.  The question
remaining is whether the observer can recover the state space itself: whether
the orbit of measurements uniquely identifies where in $X$ they are.

The answer depends on how much the delay orbit
$\{h \circ T^n : n \in \mathbb{Z}\}$ collectively distinguishes points of $X$.
If $h$ is constant, no reconstruction is possible.  At the other extreme, if
the observable $\sigma$-algebra $\mathcal{O}_h = \sigma(\{h \circ T^n\})$
equals $\mathcal{B}(X)$ modulo $\mu$, then almost every pair of distinct points
is separated by some delayed measurement.  The reconstruction theorem shows
this $\sigma$-algebraic condition is equivalent to two others: an analytic
condition on $L^2$-approximation by the delay algebra, and a geometric
condition on the delay map $\Phi_h : x \mapsto (h(T^n x))_{n \in \mathbb{Z}}$
being a measure-theoretic embedding.

The key to the equivalence is a general lemma on function algebras — the
density bridge — which says that for any subalgebra
$\mathcal{A} \subseteq L^\infty(X, \mu)$ containing $1$, the $L^2$-closure of
$\mathcal{A}$ equals $L^2(X, \sigma(\mathcal{A}), \mu)$.  The proof goes via
the functional monotone class theorem: the class of bounded measurable functions
approximable in $L^2$ by elements of $\mathcal{A}$ is a vector space closed
under bounded monotone limits, and since $\mathcal{A}$ is an algebra it already
separates points of $\sigma(\mathcal{A})$, so the class is all of
$L^\infty(\sigma(\mathcal{A}))$.  The $\sigma$-algebra and the $L^2$ space it
generates are inextricably linked.

The Koopman operator $U_T f = f \circ T$ on $L^2(X, \mu)$ provides a spectral
perspective.  The cyclic vector condition —
$\overline{\mathrm{span}}\{U_T^n h : n \in \mathbb{Z}\} = L^2$ — implies
reconstruction but is strictly stronger: it asks for density of the
\emph{linear span} of the orbit, while reconstruction asks only for density of
the \emph{algebra} the orbit generates.  For bounded $h$, the algebra contains
the linear span, so cyclicity is the stronger condition.  The two coincide
precisely when $U_T$ has simple $L^2$-spectrum.

When reconstruction holds, a compact space introduced earlier for an entirely
different purpose — the Stone space $\mathrm{St}(\mathcal{O}_h)$ of the
observable algebra, built in~\cite{mahon_paper1} as a technical completion of
the sample space for deriving $\sigma$-additivity — turns out to be
measure-theoretically isomorphic to $X$ itself.  The construction that began
by seeking a compact ambient space for the measure ends by identifying it with
the state space being observed.

Compared to the classical Takens embedding theorem~\cite{Takens1981} — which
requires a $C^2$ diffeomorphism on a compact $d$-manifold and a generic smooth
$h$, and produces a diffeomorphic embedding into $\mathbb{R}^{2d+1}$ — the
reconstruction theorem here requires only measurability, uses all finite delays
simultaneously, and replaces the dimension count with the algebraic density
condition.  On a compact manifold, Stone--Weierstrass gives density for any
point-separating $h$, recovering the measure-theoretic content of Takens
without differentiability hypotheses.

Throughout, $(X, \mathcal{B}, \mu)$ is a standard probability space and
$T : X \to X$ is a measurable, $\mu$-preserving invertible bimeasurable
bijection with $T_*\mu = \mu$.  We fix $h \in L^\infty(X, \mu)$.

\begin{definition}[Observable algebra]
\label{def:obs-algebra}
The \emph{observable algebra} generated by $h$ is
\[
  \mathcal{O}_h \;:=\; \sigma\!\bigl(\{h \circ T^n : n \in \mathbb{Z}\}\bigr)
  \;\subseteq\; \mathcal{B}(X).
\]
The \emph{delay algebra} $\mathcal{A}_h$ is the smallest subalgebra of
$L^\infty(X, \mu)$ containing $1$ and every $h \circ T^n$; concretely, finite
sums $\sum_\alpha c_\alpha \prod_k (h \circ T^{n_k})^{e_k}$.
\end{definition}

\begin{definition}[Delay map]
\label{def:delay-map}
The \emph{delay map} is
\[
  \Phi_h : X \to \mathbb{R}^\mathbb{Z},
  \qquad
  \Phi_h(x) \;=\; \bigl(h(T^n x)\bigr)_{n \in \mathbb{Z}}.
\]
Since $\pi_n \circ \Phi_h = h \circ T^n$ for each coordinate projection $\pi_n$,
the map $\Phi_h$ is measurable and
$\mathcal{O}_h = \Phi_h^{-1}(\mathcal{B}(\mathbb{R}^\mathbb{Z}))$.
Thus $\mathcal{O}_h = \mathcal{B}$ mod $\mu$ is precisely the condition that
$\Phi_h$ is a measure-theoretic embedding — the geometric face of
condition~(iii) in the reconstruction theorem below.
\end{definition}

% ====================================================================
\section{Density bridge and reconstruction}
\label{sec:density-bridge}

The central lemma connects the Boolean algebra level to the $L^2$ level.
It holds for any subalgebra of $L^\infty$.

\begin{lemma}[Density bridge]
\label{lem:density-bridge}
Let $(X, \mathcal{B}, \mu)$ be a probability space and let
$\mathcal{A} \subseteq L^\infty(X, \mu)$ be a subalgebra containing the
constant function~$1$.  Then:
\[
  \overline{\mathcal{A}}^{L^2}
  \;=\;
  L^2(X,\, \sigma(\mathcal{A}),\, \mu).
\]
Consequently, $\overline{\mathcal{A}}^{L^2} = L^2(X, \mu)$ if and only if
$\sigma(\mathcal{A}) = \mathcal{B}$ modulo $\mu$.
\end{lemma}

\begin{proof}
Write $\mathcal{M} := \sigma(\mathcal{A})$ and $V := \overline{\mathcal{A}}^{L^2}$.

\textbf{Step 1: $V \subseteq L^2(X, \mathcal{M}, \mu)$.}
Every $f \in \mathcal{A}$ is $\mathcal{A}$-measurable, hence $\mathcal{M}$-measurable.
So $\mathcal{A} \subseteq L^2(X, \mathcal{M}, \mu)$.  Since $L^2(X, \mathcal{M}, \mu)$
is a closed subspace of $L^2(X, \mathcal{B}, \mu)$, taking the closure gives
$V \subseteq L^2(X, \mathcal{M}, \mu)$.

\textbf{Step 2: $L^2(X, \mathcal{M}, \mu) \subseteq V$.}
It suffices to show that every indicator $\mathbf{1}_E$ with $E \in \mathcal{M}$
lies in $V$.  By the functional monotone class theorem~\cite{Dynkin}, the
class
\[
  \mathcal{H}
  \;:=\;
  \bigl\{f \in L^\infty(X, \mathcal{M}, \mu) : f \in V\bigr\}
\]
is a vector space containing $\mathcal{A}$ and closed under bounded monotone
pointwise limits (since if $f_n \in V$ are uniformly bounded and $f_n \to f$
pointwise a.e.\ with $f \in L^\infty$, then $f_n \to f$ in $L^2$ by dominated
convergence, so $f \in V$).  Since $\mathcal{A}$ is an algebra containing $1$
and $\sigma(\mathcal{A}) = \mathcal{M}$, the monotone class theorem gives
$\mathcal{H} = L^\infty(X, \mathcal{M}, \mu)$.  In particular every
$\mathcal{M}$-indicator $\mathbf{1}_E \in V$.

Since the $\mathcal{M}$-simple functions are dense in $L^2(X, \mathcal{M}, \mu)$
and $V$ is closed, $L^2(X, \mathcal{M}, \mu) \subseteq V$.

\textbf{Conclusion.}
Steps~1 and~2 give $V = L^2(X, \mathcal{M}, \mu)$.  The final equivalence is
immediate: $V = L^2(X, \mathcal{B}, \mu)$ iff $L^2(X, \mathcal{M}, \mu) =
L^2(X, \mathcal{B}, \mu)$ iff $\mathcal{M} = \mathcal{B}$ mod $\mu$.
\end{proof}

\begin{remark}[Role of the algebra structure]
\label{rem:algebra-necessary}
The algebra structure of $\mathcal{A}$ is essential.  The corresponding
statement for the closed \emph{linear} span $\mathcal{H}_{\mathcal{A}} =
\overline{\mathrm{span}}(\mathcal{A})$ is false in general:
$\mathcal{H}_{\mathcal{A}} = L^2$ does not imply
$\sigma(\mathcal{A}) = \mathcal{B}$.

For example, take $X = [0,1]$, $\mu = $ Lebesgue, $T(x) = 2x \bmod 1$ (the
doubling map), and $h(x) = x$.  The orbit $\{h \circ T^n\}(x) = \{2^n x \bmod 1\}$
generates $\mathcal{B}([0,1])$ as a $\sigma$-algebra (the dyadic structure separates
points).  But $\mathrm{span}\{2^n x \bmod 1 : n \geq 0\}$ is a proper subspace of
$L^2[0,1]$: it does not contain $x^2$.  So the linear span can be a proper dense
subspace even when $\sigma(\mathcal{A}) = \mathcal{B}$.

In the other direction, $h$ cyclic for $U_T$ (linear span dense) always implies
$\sigma(\mathcal{A}_h) = \mathcal{B}$ mod $\mu$ — see Corollary~\ref{cor:cyclic-implies-reconstruction}.
\end{remark}

\begin{theorem}[Reconstruction theorem]
\label{thm:reconstruction}
Let $(X, \mathcal{B}, \mu, T)$ be an invertible measure-preserving system and
$h \in L^\infty(X, \mu)$.  The following are equivalent:
\begin{enumerate}[label=(\roman*)]
  \item $\mathcal{O}_h = \mathcal{B}$ modulo $\mu$;
  \item $\mathcal{A}_h$ is dense in $L^2(X, \mu)$;
  \item $\Phi_h : X \to \mathbb{R}^\mathbb{Z}$ is a measure-theoretic embedding:
    $\Phi_h^{-1}(\mathcal{B}(\mathbb{R}^\mathbb{Z})) = \mathcal{B}$ modulo $\mu$.
\end{enumerate}
\end{theorem}

\begin{proof}
\emph{(i) $\Leftrightarrow$ (ii).}
Apply the density bridge (Lemma~\ref{lem:density-bridge}) with
$\mathcal{A} = \mathcal{A}_h$ and $\sigma(\mathcal{A}_h) = \mathcal{O}_h$:
$\overline{\mathcal{A}_h}^{L^2} = L^2(X, \mu)$ iff $\mathcal{O}_h = \mathcal{B}$
mod $\mu$.

\emph{(i) $\Leftrightarrow$ (iii).}
By Definition~\ref{def:delay-map},
$\mathcal{O}_h = \Phi_h^{-1}(\mathcal{B}(\mathbb{R}^\mathbb{Z}))$.
So $\mathcal{O}_h = \mathcal{B}$ mod $\mu$ iff
$\Phi_h^{-1}(\mathcal{B}(\mathbb{R}^\mathbb{Z})) = \mathcal{B}$ mod $\mu$,
which is exactly condition~(iii).
\end{proof}

\begin{remark}[What ``measure-theoretic embedding'' means]
\label{rem:embedding}
Condition~(iii) says that $\Phi_h$ is \emph{almost everywhere injective} and
the pullback $\sigma$-algebra is all of $\mathcal{B}$.  More precisely: since
$\Phi_h^{-1}(\mathcal{B}(\mathbb{R}^\mathbb{Z})) = \mathcal{B}$ mod $\mu$, for
any two sets $A, A' \in \mathcal{B}$ with $\mu(A \triangle A') > 0$ there exists
a Borel set $E \subseteq \mathbb{R}^\mathbb{Z}$ with
$\mu(\Phi_h^{-1}(E) \triangle A) = 0$ and
$\mu(\Phi_h^{-1}(E) \triangle A') > 0$.  In measure-theoretic language,
$\Phi_h$ induces an isomorphism of measure algebras
$(X/\mu, \mathcal{B}/\mu) \cong (\mathbb{R}^\mathbb{Z}/(\Phi_h)_*\mu,
\mathcal{B}(\mathbb{R}^\mathbb{Z})/(\Phi_h)_*\mu)$.

The $T$-invariance of $\mu$ implies that $(\Phi_h)_*\mu$ is invariant under
the bilateral shift $\sigma$ on $\mathbb{R}^\mathbb{Z}$, since
$(\Phi_h(Tx))_n = h(T^{n+1}x) = (\sigma(\Phi_h(x)))_n$ for each $n$.  The
dynamical system $(X, T, \mu)$ is therefore isomorphic to
$(\mathbb{R}^\mathbb{Z}, \sigma, (\Phi_h)_*\mu)$ restricted to the image of
$\Phi_h$.
\end{remark}

\begin{definition}[Cyclic vector]
\label{def:cyclic}
The function $h \in L^2(X, \mu)$ is a \emph{cyclic vector} for $U_T$ if
\[
  \mathcal{H}_h
  \;:=\;
  \overline{\mathrm{span}}\!\bigl\{U_T^n h : n \in \mathbb{Z}\bigr\}
  \;=\; L^2(X, \mu).
\]
\end{definition}

\begin{corollary}[Cyclic implies reconstruction]
\label{cor:cyclic-implies-reconstruction}
If $h \in L^\infty(X, \mu)$ is a cyclic vector for $U_T$, then
$\mathcal{O}_h = \mathcal{B}$ modulo $\mu$.
\end{corollary}

\begin{proof}
Since $\mathcal{A}_h$ contains $\mathrm{span}\{h \circ T^n\} =
\mathrm{span}\{U_T^n h\}$ as a subset, $\overline{\mathcal{A}_h}^{L^2}
\supseteq \mathcal{H}_h = L^2$.  By the density bridge, $\mathcal{O}_h =
\mathcal{B}$ mod $\mu$.
\end{proof}

\begin{remark}[The implication is strict]
\label{rem:strict}
The converse of Corollary~\ref{cor:cyclic-implies-reconstruction} fails in
general: $\mathcal{O}_h = \mathcal{B}$ does not imply $h$ is cyclic.
The doubling map example of Remark~\ref{rem:algebra-necessary} witnesses this.

The cyclic vector condition asks for density of the \emph{linear span} of the
orbit; reconstruction asks only for density of the \emph{algebra}.  For a bounded
$h$, the linear span is contained in the algebra, so cyclicity is the stronger
condition.

The two conditions coincide precisely when $U_T$ has \emph{simple
$L^2$-spectrum}: $L^2(X,\mu)$ is a cyclic $U_T$-module, meaning there exists
some cyclic vector.  In this case, every generating observation $h$ (satisfying
$\mathcal{O}_h = \mathcal{B}$) is automatically cyclic.  Simple spectrum is
a generic property of ergodic transformations in the Baire category sense
\cite{Halmos1944} but is not universal.
\end{remark}

The Stone space $\mathrm{St}(\mathcal{O}_h)$ of the observable algebra is the
compact Hausdorff space whose points are ultrafilters on $\mathcal{O}_h$, with
the natural Stone embedding $\iota : X \to \mathrm{St}(\mathcal{O}_h)$ sending
$x$ to the principal ultrafilter $\{E \in \mathcal{O}_h : x \in E\}$.

\begin{proposition}[Stone space identification]
\label{prop:stone-id}
Suppose $\mathcal{O}_h = \mathcal{B}$ modulo $\mu$.  Then $\iota$ is a
measure-theoretic isomorphism: $\iota$ identifies the measure algebra
$(X/\mu, \mathcal{B}/\mu)$ with
$(\mathrm{St}(\mathcal{O}_h)/\iota_*\mu,
\mathcal{B}(\mathrm{St}(\mathcal{O}_h))/\iota_*\mu)$.
\end{proposition}

\begin{proof}
\emph{Injectivity mod $\mu$.}  Suppose $\iota(x) = \iota(x')$.  For any
$E \in \mathcal{B}$ with $x \in E \not\ni x'$, find $F \in \mathcal{O}_h$
with $\mu(E \triangle F) = 0$.  Since $x \in F$ and $\iota(x) = \iota(x')$,
also $x' \in F$, so $x' \in (E \triangle F) \cup F^c$, placing $x'$ in a
null set.  Hence $\iota$ is injective $\mu$-a.e.

\emph{$\sigma$-algebra identification.}  The pullback $\iota^{-1}$ maps the
Baire $\sigma$-algebra of $\mathrm{St}(\mathcal{O}_h)$ (generated by the
clopen sets $[E] = \{u : E \in u\}$) to $\mathcal{O}_h = \mathcal{B}$ mod
$\mu$.

\emph{Measure recovery.}  Since $\mu$ is $\sigma$-additive on
$\mathcal{O}_h = \mathcal{B}$, the collective exhaustion condition
of~\cite{mahon_paper1} forces $\iota_*\mu$ to concentrate on the principal
ultrafilters $\iota(X)$, so $\iota_*\mu$ is supported on the image and the
measure algebras are identified.
\end{proof}

\noindent The Stone space introduced in~\cite{mahon_paper1} as a compact
ambient space for deriving $\sigma$-additivity is, under the reconstruction
condition, the state space $X$ itself.  The construction that began by seeking
a compact completion of the sample space ends by recognising it as the object
the observer was trying to recover.

The classical Takens embedding theorem~\cite{Takens1981} reaches a related
conclusion by a different route: for a $C^2$ diffeomorphism on a compact
$d$-manifold and a generic smooth $h$, the finite-delay map
$\Phi_h^{(k)}(x) = (h(x), h(Tx), \ldots, h(T^{k-1}x))$ is a diffeomorphic
embedding into $\mathbb{R}^{2d+1}$ for $k \geq 2d+1$.  The reconstruction
theorem here requires only measurability, uses all delays simultaneously, and
replaces the dimension count with the algebraic density condition; in exchange,
the embedding is measure-theoretic rather than topological.  On a compact
manifold with generic smooth $h$, the orbit $\{h \circ T^n\}$ separates
points, so Stone--Weierstrass gives $\mathcal{A}_h$ uniformly dense in $C(M)$
and hence $L^2$-dense: the density condition holds and the two results are
compatible.

The question the theorem leaves open is the genericity of the reconstruction
condition in the purely measure-theoretic setting.  On compact manifolds the
analysis is classical; on abstract standard probability spaces, characterising
the pairs $(T, h)$ for which $\mathcal{O}_h = \mathcal{B}$ modulo $\mu$ — and
understanding how the answer varies with the spectral properties of $T$ — is
the natural next problem.

% ====================================================================
\bibliographystyle{plain}
\bibliography{references}
